\documentclass[10pt,a4paper]{amsart}
\usepackage[margin=19mm]{geometry}
\usepackage{amsmath,amssymb,mathtools}
\usepackage[scr=rsfs,cal=euler]{mathalfa}
\usepackage{xfrac}
\usepackage[unicode,hidelinks,bookmarks=false]{hyperref}
\newcommand{\ie}{i.e.\ }
\newcommand{\cf}{cf.\ }
\newcommand{\resp}{respectively\ }
\newcommand{\Subsection}{\subsection}
\providecommand{\nobreakdash}{\nobreak\hbox{-}}
\setlength{\emergencystretch}{2em}
\multlinegap0pt
\usepackage{amssymb}
\newtheorem{lemma}{Lemma}
\newtheorem{theorem}{Theorem}
\title{On potentials of distributions in Orlicz-Hardy type spaces on the Heisenberg group}
\author{Pablo Rocha}
\date{Source: arXiv:2605.24194v1 (CC BY 4.0)}
\begin{document}
\maketitle

\noindent\emph{Source and licence.} On potentials of distributions in Orlicz-Hardy type spaces on the Heisenberg group. Version arXiv:2605.24194v1 (CC BY 4.0), distributed by arXiv under the Creative Commons Attribution 4.0 International licence, http://creativecommons.org/licenses/by/4.0/, as recorded in the arXiv licence metadata for this exact version and shown on the abstract page https://arxiv.org/abs/2605.24194v1. Full licence notice: \texttt{licenses/arxiv-2605.24194-CC-BY-4.0.txt}.\par\smallskip

\noindent\textit{Editorial scope.} Counted unit: the article's theorem thm (source0.tex lines 1282--1285). The article's complete original proof of it is delivered with it (source0.tex lines 1287--1309, one \texttt{proof} environment of 23 lines). No mathematics is written, repaired or re-derived: the statement and the proof are the article's own lines, in the article's own notation, and no retained line is substituted or re-flowed.  Retained uncounted as the source-local context the proof needs: lines 387--393; lines 405--408.  Nothing was omitted from any retained span, so the package carries no omission record.  No retained line contains a citation, so the wrapper carries no bibliography and no bibliography entry.  No retained line contains a figure environment, an image-inclusion command or drawing code, so no image is delivered and the package declares no asset dependency.  Editorial formatting only, in addition: an AMS \texttt{amsart} preamble that restates the article's own declarations and macros used by the retained lines, namely the environments \texttt{lemma}, \texttt{theorem} on separate counters, together with the standard packages those lines need.  The retained environments print in this document as Lemma 1, Theorem 1 in order of appearance, whereas the article prints the same items with the numbering of its own full document.  The complete native source of the article as distributed in the arXiv e-print for this exact version is bundled unchanged as \texttt{sources/201258/source0.tex}.
\begin{lemma} \label{lower estim}
Let $\Phi$ be an Orlicz function with positive lower type $p_{\Phi}^{-}$ and positive upper type $p_{\Phi}^{+}$. Then,

(i) $\Phi(t) \gtrsim \, t^{p_{\Phi}^{-}}$, if \, $t \geq 1$,

(ii) $\Phi(t) \gtrsim \, t^{p_{\Phi}^{+}}$, if \, $0 < t \leq 1$.
\end{lemma}

\begin{lemma} \label{modular 1}
Let $\Phi$ be an Orlicz function of positive upper type $p_{\Phi}^{+}$. Then, $f \in L^{\Phi}(\mathbb{H}^n)$ if and only if 
$\kappa_{\Phi}(f) < \infty$.
\end{lemma}

\begin{theorem} \label{2nd thm}
If $1 < q < \frac{n+1}{n}$ and $\Phi$ is an Orlicz function with $0 < I(\Phi) < Q \, (2 + \frac{Q}{q})^{-1}$, then 
$\mathcal{H}^{\Phi}_{q, \, 2}(\mathbb{H}^{n}) = \{ 0 \}.$
\end{theorem}

\begin{proof} Let $G \in \mathcal{H}^{\Phi}_{q, \, 2}(\mathbb{H}^{n})$ and assume $G \neq 0$. Then there exists $g \in G$ 
that is not a polynomial of homogeneous degree less or equal to $1$. It is easy to check that there exist a positive constant 
$c \in (0, 1]$ and a $\rho$ - ball $B = B(e, r)$ with $r > 1$ such that
\[
\int_{B} |g(w) - P(w)|^{q} \, dw \geq c > 0,
\]
for every $P \in \mathcal{P}_{1}$.

Let $z$ be a point such that $\rho(z) > r$ and let $\delta = 2 \rho(z)$. Then $B(e, r) \subset B(z, \delta)$. If $h \in G$, then 
$h = g - P$ for some $P \in \mathcal{P}_{1}$ and
\[
\delta^{-2}|h|_{q, B(z, \delta)} \geq c \rho(z)^{-2-Q/q}.
\]
So $N_{q,2}(G; \, z) \geq c \, \rho(z)^{-(2+Q/q)}$, for $\rho(z) > r$. Since $\Phi$ is an Orlicz function of positive upper type 
$p_{\Phi}^{+}$ with $I(\Phi) \leq p_{\Phi}^{+} < Q(2+Q/q)^{-1}$ and $0 < c \, \rho(z)^{-(2+Q/q)} < 1$ on $\rho(z) > r$, by 
Lemma \ref{lower estim} - (ii), we have that
\begin{eqnarray*}
\int_{\mathbb{H}^n} \Phi\left( N_{q,2}(G; z) \right) dz &\geq& \, \int_{\rho(z) > r} \Phi\left( c \rho(z)^{-(2+Q/q)} \right) \, dz \\
&\gtrsim& \int_{\rho(z) > r} \rho(z)^{-(2+Q/q)p^{+}_{\Phi}} \, dz = \infty,
\end{eqnarray*}
which gives a contradiction (see Lemma \ref{modular 1}). Thus $\mathcal{H}^{\Phi}_{q, 2}(\mathbb{H}^{n}) = \{0\}$, 
if $0 < I(\Phi) < Q(2+Q/q)^{-1}$.
\end{proof}

\end{document}
